This appendix records a separate compression result for the retained
universal compiler. It is not offered as a better fixed Con(ZFC)
equation than the compact forms in Section~\ref{sec:quotient}.
The complete source and ordinary proof are under
\file{additional/ordinate_normalization/}.

Use the parent notation $\Delta=(a+2)^2-1$, with $c,i,T>0$ and
integer rank $R$, and put
\[
 I=ic^2,\quad S=\Delta I,\quad Q=S^2,\quad D=1+IS,
 \quad u=cT,\quad b=c+RD,\quad p=Du^2.
\]
For a new positive coordinate $z$, define
\[
 q_z=Qz,\qquad W=1+(q_z-z)(q_z-2),\qquad
 A=p+b^2-W,\qquad K=A^2-4pb^2.
\]
If $P_5$ is the product of the retained first, main, input, index and
transport factors, the new equation is
\[
 P_5(K+1)-1=0.
\]
It preserves the ordinary input of the universal compiler, with 17
positive witnesses, 91 operations and exact degree 269. The matched
older root-elimination construction has 92 operations and degree 325.
This universal count separates the ordinary input from the witnesses.

The mechanism is to impose $y=Qz-1$ before eliminating the strong
square root $f$. Indeed
\[
 QV^2-(Q-1)y^2-1=Q(V^2-W).
\]
The nonzero factor $Q$ cancels before taking the quadratic norm.
This avoids squaring a larger coefficient in the elimination step.

For soundness, the positive domain gives $Q>1$, $D>c$, $u>0$ and
$b\ne0$, even before a bound on $R$ is available. At a zero, the
integer unit $K+1$ is 1 or 2 modulo 4, so it is $+1$; hence
$P_5=1$ and $K=0$. The proof-only rational value
\[
 f=\frac{A}{2ub}
\]
satisfies $f^2=D$. A rational square root of an integer is integral.
Set $V=uf-b$ and $y=Qz-1>0$. Then $V^2=W$, the old auxiliary norm
is one, and the normalized strong norm is one.

The inherited source-specific signed-quotient theorem is needed to
recover $f>0$. If $f<0$, replacing $(f,T)$ by $(-f,-T)$ leaves all
old products unchanged and produces the forbidden positive-$f$,
negative-$T$ zero. This is an explicitly retained theorem for the valid
compiler recipe; it is not asserted for arbitrary numerical ports.

For completeness the canonical parent uses $R\equiv3\pmod4$ and
$y_* =\psi_S(R)$. Equation~\eqref{eq:oddres} gives $Q\mid y_*+1$,
so $z=(y_*+1)/Q>0$ is integral. The canonical strong and auxiliary
completion then makes $K=0$. This preserves accepted-input existence,
not every old coordinate tuple.

The degree calculation uses $\deg P_5=81$, $\deg Q=46$,
$\deg W=94$, $\deg A=94$ and $\deg K=188$.
The full nonzero leading form and an attained coefficient
$32\,\mathrm{Bm1}^{162}$ are supplied. Finite source and algebra
checks accompany the ordinary proof and its stated parent interface.
